\documentclass[UTF8,fontset=windows,zihao=-4,a4paper]{ctexart}
\usepackage[left=24mm,right=24mm,top=23mm,bottom=23mm]{geometry}
\usepackage{xcolor,fancyhdr,hyperref,needspace}
\definecolor{ink}{HTML}{173D35}
\hypersetup{colorlinks=true,linkcolor=ink,pdftitle={校园活动工作台汇报讲稿},pdfauthor={张敬宇、刘秋彤、王梓安、胡辰凯}}
\pagestyle{fancy}\fancyhf{}
\fancyhead[L]{\small 校园活动工作台汇报讲稿}\fancyhead[R]{\small 5 分钟主讲}
\fancyfoot[C]{\small\thepage}\setlength{\headheight}{15pt}
\setlength{\parindent}{0pt}\setlength{\parskip}{5pt}\linespread{1.22}
\begin{document}
\begin{center}{\Large\bfseries\color{ink}校园活动工作台汇报讲稿}\end{center}
对应当前 17 页展示稿。第 1–16 页为主讲，建议约 5 分钟；第 17 页为附录，提问时再使用。按自然语速讲述，页码和时间提示不用读出。

\Needspace{8\baselineskip}
{\large\bfseries\color{ink}第 1 页：校园活动，从策划到收尾}\par
{\small\color{gray}建议 10 秒，0:00–0:10}\par
大家好，我们小组制作的是校园活动策划与管控工作台。下面沿着一场活动的组织过程，介绍它能完成哪些工作。\par\vspace{8pt}

\Needspace{8\baselineskip}
{\large\bfseries\color{ink}第 2 页：校园活动组织的背景与问题}\par
{\small\color{gray}建议 15 秒，0:10–0:25}\par
组织校园活动时，计划、人员和材料需要不断协调。比如场地调整，就会影响宣传、设备和应急安排。我们希望用同一个工作台，让信息保持一致，也让责任和交付衔接起来。\par\vspace{8pt}

\Needspace{8\baselineskip}
{\large\bfseries\color{ink}第 3 页：汇报目录}\par
{\small\color{gray}建议 10 秒，0:25–0:35}\par
汇报分为四部分：背景与设计、核心功能、风险与交付，以及结构化方法和运行验证。先看整体功能，再按使用顺序展开。\par\vspace{8pt}

\Needspace{8\baselineskip}
{\large\bfseries\color{ink}第 4 页：把活动组织的各项工作连接起来}\par
{\small\color{gray}建议 15 秒，0:35–0:50}\par
工作台以项目为中心，连接任务、人员资源、AI材料、报名考勤、风险和成果。组织者统筹项目，成员完成任务，参与者通过手机填写，各自的操作在同一个活动中关联。\par\vspace{8pt}

\Needspace{8\baselineskip}
{\large\bfseries\color{ink}第 5 页：一条完整的活动组织路径}\par
{\small\color{gray}建议 10 秒，0:50–1:00}\par
活动分为策划、筹备、执行和收尾四阶段。每阶段都对应操作和材料，协作与风险管理贯穿全过程，最后还能归档复用。\par\vspace{8pt}

\Needspace{8\baselineskip}
{\large\bfseries\color{ink}第 6 页：先创建项目，再统一活动事实}\par
{\small\color{gray}建议 20 秒，1:00–1:20}\par
登录后先进入个人项目库，可以新建、导入或复制活动。进入项目后，日期、场地、人数和预算集中维护，后续任务和材料沿用这些信息。策划时还可以比较不同方案，保留评分和选择理由。\par\vspace{8pt}

\Needspace{8\baselineskip}
{\large\bfseries\color{ink}第 7 页：把“活动计划”变成可追踪的任务}\par
{\small\color{gray}建议 20 秒，1:20–1:40}\par
任务页把计划落实为负责人、期限、前置依赖和交付成果。成员按权限协作，任务从未开始推进到已验收。共享保存会检查版本，发生冲突时保留草稿，让执行和交接都有明确依据。\par\vspace{8pt}

\Needspace{8\baselineskip}
{\large\bfseries\color{ink}第 8 页：排好人员班次，同时发现资源缺口}\par
{\small\color{gray}建议 20 秒，1:40–2:00}\par
组织者可以为多人安排独立班次，系统检查时间冲突、技能要求和资源占用。资源台账统一核对库存与需求。图中提示麦克风和运输车辆不足，缺口没有解决前，系统会阻止整组安排保存。\par\vspace{8pt}

\Needspace{8\baselineskip}
{\large\bfseries\color{ink}第 9 页：从策划书到宣传海报，形成可用材料}\par
{\small\color{gray}建议 30 秒，2:00–2:30}\par
材料模块支持二十一类产出，覆盖策划书、预算、通知、执行手册和复盘决算等。文字使用DeepSeek，海报使用GPT-image2。这里展示已有策划书和实际生成的宣传海报。生成后可以在线修改、补齐表格，再由使用者核对事实，确认交付版本，导出Word、Excel或Markdown文件。\par\vspace{8pt}

\Needspace{8\baselineskip}
{\large\bfseries\color{ink}第 10 页：让资料和上游结果进入下一项工作}\par
{\small\color{gray}建议 20 秒，2:30–2:50}\par
资料库保存参考内容和来源，由使用者选择哪些进入AI输入。卡片工作流关联项目事实、资料和上游结果，例如先统一活动安排，再生成参与者通知和岗位简报。提示词明确目标、依据和输出要求，使材料保持一致。\par\vspace{8pt}

\Needspace{8\baselineskip}
{\large\bfseries\color{ink}第 11 页：活动现场：手机填写，按班次执行}\par
{\small\color{gray}建议 30 秒，2:50–3:20}\par
进入现场环节，参与者可以通过链接或二维码在手机上报名，提交的信息汇总到项目。工作人员则在个人任务页查看班次，接受安排、签到签退，并填写执行报告。两类人员分别管理，组织者可以核对到场和执行情况。现场记录还会进入后续交付与复盘，让安排和实际行动对应起来。\par\vspace{8pt}

\Needspace{8\baselineskip}
{\large\bfseries\color{ink}第 12 页：风险有评分，也有对应的应急动作}\par
{\small\color{gray}建议 30 秒，3:20–3:50}\par
风险台账用可能性乘影响评分，并配套预防、响应和升级措施。以户外活动遇到恶劣天气为例，评分十六，属于高风险。预案需要落实备用场地、决策时间、通知和转场责任。工作台还支持场景演练，逐步记录行动，结束后把结论写回风险台账，使预案能够根据检查结果修订。\par\vspace{8pt}

\Needspace{8\baselineskip}
{\large\bfseries\color{ink}第 13 页：交付可验收，收尾可归档与复用}\par
{\small\color{gray}建议 20 秒，3:50–4:10}\par
任务完成后，责任人提交已确认的材料版本，组织者验收或退回。后续编辑会保留历史验收依据。收尾时继续核对费用与物资，形成决算、复盘和归档文件，再导出项目包，为下一场活动复用。\par\vspace{8pt}

\Needspace{8\baselineskip}
{\large\bfseries\color{ink}第 14 页：结构化方法怎样进入实际操作}\par
{\small\color{gray}建议 20 秒，4:10–4:30}\par
结构化方法落实在实际操作中：SPIN明确需求，MECE和逻辑树分解任务，5W2H补齐责任与资源。ALIGN和RTGO规定AI的依据与输出，PDCA让验收、演练和修订进入下一轮。右侧方案比选也使用统一维度和权重。\par\vspace{8pt}

\Needspace{8\baselineskip}
{\large\bfseries\color{ink}第 15 页：运行协作、测试与小组贡献}\par
{\small\color{gray}建议 20 秒，4:30–4:50}\par
不同设备通过可达地址和邀请权限协作。离线回归覆盖三十八个模块、一千零四十四项断言，全部通过。分工上，张敬宇负责整体设计与功能制作；刘秋彤负责界面交互、功能完善和展示材料整合，两人共同完成报告与PPT。王梓安负责测试，胡辰凯负责流程、风险预案和材料核对。\par\vspace{8pt}

\Needspace{8\baselineskip}
{\large\bfseries\color{ink}第 16 页：我们交付的是一套完整的组织工具}\par
{\small\color{gray}建议 10 秒，4:50–5:00}\par
我们的工作台覆盖全周期，形成过程材料，并配套分级应急预案。AI辅助生成，使用者负责判断、修订和验收。谢谢大家。\par\vspace{8pt}

\Needspace{8\baselineskip}
{\large\bfseries\color{ink}第 17 页：功能与中间产物索引}\par
{\small\color{gray}附录，不计入主讲时长}\par
这一页按功能和活动阶段列出操作空间与产出类型，可以在提问时查阅。二十一类指支持的材料类型，项目可按需启用，也能增加自定义产出。\par\vspace{8pt}
\end{document}